% =============================================================================
% Kapitel 6: Implementierung
% =============================================================================

Dieses Kapitel beschreibt die technische Umsetzung des in Kap.~\ref{sec:entwurf} entworfenen Systems.
Zunächst wird die Framework-Migration behandelt, anschließend die Integration der Embedding-basierten Analysefunktionen und abschließend die Gestaltung der Benutzeroberfläche.

\subsection{Framework-Migration}\label{sec:migration}

Die in Kap.~\ref{sec:vorarbeit} identifizierten Defizite lassen sich in zwei Kategorien einteilen: technische Defizite, die durch die Framework-Migration adressiert werden, und funktionale Lücken, die durch die Embedding-basierte Erweiterung geschlossen werden sollen.
Dieser Abschnitt beschreibt die Migrationsentscheidungen; die Embedding-Integration folgt in Kap.~\ref{sec:impl-embedding}.

\subsubsection{Technologieentscheidungen}\label{sec:tech-entscheidungen}

Die fehlende Supplier Stability (NFA4) erfordert die Migration der drei zentralen Abhängigkeiten: Nuxt~2 wird durch Nuxt~4, Vue~2 durch Vue~3 und Vuetify~2 durch Vuetify~4 ersetzt.
Damit werden ausschließlich aktiv gewartete Frameworks eingesetzt, die Security-Patches und langfristigen Support erhalten.

Der in der Anwendungsbewertung kritisierte Mangel an Separation of Concerns -- die gesamte Anwendungslogik in einer 1334-zeiligen Datei -- wird durch den Wechsel von der Options API auf die Composition API adressiert.
Vue~3 erlaubt es, zusammengehörige Logik in sogenannte Composables zu extrahieren, anstatt sie nach Optionstyp (\texttt{data}, \texttt{methods}, \texttt{computed}) zu gruppieren.
Dadurch wird die monolithische Editorkomponente in spezialisierte Module aufgeteilt.

Die von Loth und Konert selbst identifizierte Limitation~L1 -- die fehlerhafte Texthervorhebung im Volltexteditor~\cite{loth2022kompetenzeditor} -- wird durch den Ersatz der ContentEditable-Komponente durch Tiptap~3 adressiert.
Der Prototyp verwendete die Browser-API \texttt{ContentEditable} direkt, was dazu führte, dass HTML-Manipulation und Cursorposition manuell synchronisiert werden mussten.
Tiptap basiert auf ProseMirror und arbeitet mit einem strukturierten Dokumentenmodell, in dem Textmarkierungen deklarativ als sogenannte Marks definiert werden.
Die Farbcodierung der Verben wird über Custom Marks realisiert, die den vier Kategorien des Prototyps entsprechen: empfohlen, nicht empfohlen, NLP-erkannt und nicht zugeordnet.

Das inkonsistente Configuration Management wird durch den Wechsel von Vuex auf Pinia sowie durch den Einsatz auto-importierter Composables anstelle globaler Plugins vereinheitlicht.
Das informelle Datenmodell -- im Prototyp als hardcodierte Seed-Daten in einer Plugin-Datei realisiert -- wird in ein separates Datenmodul überführt, das ein explizites Schema definiert.
Die Datenbankschicht nutzt Dexie~4 und kapselt alle Verblisten, Taxonomie-Zuordnungen und CRUD-Operationen.

\subsubsection{Architektur der migrierten Anwendung}\label{sec:architektur-migration}

Kern der Restrukturierung ist die Aufteilung der Editorkomponente in vier Composables, die jeweils eine abgegrenzte Verantwortlichkeit übernehmen.
Das Composable \texttt{useDatabase} kapselt die Dexie-Anbindung und stellt Verblisten, Persistenz und CRUD-Operationen bereit; es adressiert den Befund des informellen Datenmodells, indem das Schema explizit definiert wird.
\texttt{useVerbAnalysis} enthält die Kernlogik der regelbasierten Analyse: Verb-Extraktion, Abgleich mit der Verbliste und Farbzuordnung -- Funktionalität, die im Prototyp über mehrere Methoden verteilt war.
\texttt{useSpacy} isoliert die Kommunikation mit dem spaCy-Backend für POS-Tagging als eigenständige Abhängigkeit, sodass diese Schnittstelle später durch das Embedding-Backend ersetzt werden kann.
\texttt{useScoring} extrahiert die Berechnung des Qualitäts-Scores, die im Prototyp in die UI-Komponente eingebettet war.

Die Seitenkomponente selbst beschränkt sich auf Template und UI-Steuerung; die fachliche Logik wird vollständig über die Composables bezogen.

\subsubsection{Erhaltene Funktionalität}\label{sec:erhaltene-funktionalitaet}

Die Migration verfolgt das Ziel, die in Tab.~\ref{tab:ist-soll} als erfüllt bewerteten Anforderungen zu erhalten.
Das Echtzeit-Feedback (FA5) wird durch denselben Debounce-Mechanismus realisiert -- nach einer Tipp-Pause wird die Analyse angestoßen und das Ergebnis als farbliche Hervorhebung im Editor dargestellt.
Die Transparenz (NFA2) bleibt durch die regelbasierte Zuordnung gewährleistet, die deutsche Sprachunterstützung (NFA3) durch die unveränderte Verbliste und das spaCy-Modell sichergestellt, und die Performance (NFA5) durch die weiterhin clientseitige regelbasierte Analyse gegeben.

Die Persistenz (FA6) wird gegenüber dem Prototyp verbessert: Die Speicherfunktion unterscheidet nun zwischen Neuanlage und Aktualisierung bestehender Einträge und adressiert damit den in der IST-Analyse dokumentierten Fehler, bei dem jeder Speichervorgang einen neuen Eintrag erstellte.

\subsubsection{Reproduzierbarkeit und Orchestrierung}\label{sec:orchestrierung}

Die IST-Analyse identifizierte zwei zusammenhängende Defizite: Das spaCy-Backend fehlt im Repository, und das System ist nicht out-of-the-box reproduzierbar (vgl. Kap.~\ref{sec:technische-bewertung}, Failure Rate).
Beide Defizite werden durch eine Docker-Compose-Konfiguration adressiert, die das spaCy-Backend als Container bereitstellt.
Ein einzelner Befehl startet sowohl das NLP-Backend als auch den Entwicklungsserver der Anwendung; die manuelle Einrichtung separater Dienste entfällt damit.
Die Nuxt-Konfiguration leitet API-Anfragen über einen Proxy an den spaCy-Container weiter, sodass Frontend und Backend ohne manuelle URL-Konfiguration zusammenarbeiten.

Das in der IST-Analyse festgestellte Fehlen von CI/CD und automatisierten Tests (vgl. Kap.~\ref{sec:technische-bewertung}, Configuration Management) wird im Rahmen dieser Arbeit nicht vollständig adressiert, da der Fokus auf der Embedding-basierten Analysefunktionalität liegt.
Die Architekturentscheidung für isolierte Composables schafft jedoch die Voraussetzung für eine spätere Testbarkeit: Jedes Composable kann unabhängig von der UI-Schicht getestet werden.

\subsection{Embedding-Integration}\label{sec:impl-embedding}

Die in Kap.~\ref{sec:entwurf} entworfene Hybrid-Pipeline wird als Python-Backend realisiert, das drei Dienste kapselt: die Analyse-Pipeline, den Embedding-Service und die Ähnlichkeitssuche.
Dieser Abschnitt beschreibt die Implementierung dieser Komponenten.

\subsubsection{Backend-Service}\label{sec:backend-service}

Der Backend-Service nutzt FastAPI mit dem Lifespan-Pattern: Beim Serverstart werden alle Modelle und Daten synchron in den Speicher geladen, bevor der erste Request bearbeitet wird.
Das spaCy-Modell \texttt{de\_core\_news\_lg}, das SBERT-Modell, die beiden SVM-Klassifikatoren und die Referenz-Embeddings stehen anschließend für die gesamte Laufzeit im Speicher bereit.

Die API ist bewusst auf zwei Endpunkte beschränkt.
Der Health-Endpunkt (\texttt{GET /health}) gibt den Ladestatus der Modelle zurück, sodass das Frontend prüfen kann, ob das Backend betriebsbereit ist.
Der Analyse-Endpunkt (\texttt{POST /analyze}) nimmt den Eingabetext entgegen, delegiert die Verarbeitung an die Pipeline und gibt das strukturierte Ergebnis zurück.
CORS ist für die Entwicklungsumgebung offen konfiguriert; die Nuxt-Konfiguration leitet API-Anfragen über einen Proxy an den Backend-Service weiter, sodass Frontend und Backend unter derselben Origin operieren.

\subsubsection{Analyse-Pipeline}\label{sec:impl-pipeline}

Die Klasse \texttt{AnalysisPipeline} orchestriert die sechs Verarbeitungsschritte aus Kap.~\ref{sec:cascading-entwurf}.
Sie kapselt die Abhängigkeiten zu den drei Diensten (spaCy, Embedding-Service, Similarity-Search) und definiert interne Datenstrukturen für Zwischenergebnisse.

Für jeden Satz wird zunächst die linguistische Vorverarbeitung durchgeführt: spaCy extrahiert Verben mit POS-Tag und Lemma und prüft über Dependency-Parsing, ob \emph{die Studierenden} als Subjekt auftreten.
Die regelbasierte Klassifikation operiert auf den extrahierten Lemmata und gibt entweder ein eindeutiges Ergebnis zurück oder \texttt{None}, was den Embedding-Fallback auslöst.

Für Effizienz werden die Embeddings aller Sätze in einem Batch berechnet, anstatt einzeln.
Der SBERT-Encoder verarbeitet Listen von Sätzen mit einer konfigurierbaren Batch-Size (Standard: 64), was die GPU-Auslastung bei mehreren Sätzen verbessert.

Die Set-Analyse wird nach der Einzelsatz-Analyse berechnet: Die Taxonomiestufen aller als Kompetenzformulierung klassifizierten Sätze werden aggregiert, die Abdeckung (Anzahl belegter Stufen von sechs) ermittelt und eine textuelle Empfehlung generiert, die auf fehlende Stufen hinweist.

\subsubsection{Embedding-Service}\label{sec:impl-embedding-service}

Der \texttt{EmbeddingService} kapselt das SBERT-Modell und die beiden SVM-Klassifikatoren.
Das SBERT-Modell wird über die sentence-transformers-Bibliothek geladen; die SVMs werden als serialisierte sklearn-Pipelines (\texttt{.pkl}-Dateien) eingelesen.

Die Encoding-Methode überführt Sätze in 768-dimensionale Vektoren.
Für die Ähnlichkeitssuche steht eine zusätzliche Methode bereit, die die Vektoren auf Einheitslänge normalisiert, sodass die Cosine-Similarity durch ein einfaches Dot-Produkt berechnet werden kann.

Die Klassifikationsmethoden geben neben der Vorhersage einen Konfidenzwert aus, der aus \texttt{predict\_proba()} der SVM abgeleitet wird -- die maximale Klassenwahrscheinlichkeit über die Platt-Skalierung.
Dieser Wert wird im Frontend als Prozentangabe angezeigt und adressiert die Anforderung NFA1 (Konfidenzmaße).

\subsubsection{Ähnlichkeitssuche}\label{sec:impl-similarity}

Die Klasse \texttt{SimilaritySearch} lädt beim Start die vorberechneten Referenz-Embeddings als NumPy-Matrix ($n \times 768$, wobei $n$ der Anzahl der Referenzsätze im Gold-Standard entspricht) und die zugehörigen Metadaten.
Die \texttt{find\_similar()}-Methode berechnet die Cosine-Similarity eines Eingabe-Embeddings gegen alle Referenzen durch eine einzelne Matrix-Vektor-Multiplikation.

Vor der Rückgabe werden Filter angewandt: Der Eingabesatz selbst wird über Textvergleich ausgeschlossen, und optional kann nach Satztyp oder Taxonomiestufe gefiltert werden.
Die Methode gibt die Top-$k$ Treffer mit Similarity-Wert, Satztext, Taxonomiestufe und Modulherkunft zurück; Treffer mit Similarity $\leq$ 0 werden verworfen.

\subsection{Benutzeroberfläche}\label{sec:impl-ui}

Die Benutzeroberfläche setzt das dreispaltige Layout des Prototyps in aktueller Technologie um und erweitert es um die Darstellung der Embedding-basierten Analyseergebnisse.
Das Frontend ist als Nuxt-4-Anwendung mit Vue~3 Composition API und Vuetify~4 implementiert.

\subsubsection{Editor und Verb-Highlighting}\label{sec:impl-editor}

Der Editor basiert auf Tiptap~3, einem ProseMirror-Wrapper, der ein strukturiertes Dokumentenmodell bereitstellt.
Die Verb-Hervorhebung wird über eine Custom Extension realisiert, die ProseMirror-Decorations verwendet.
Decorations sind Inline-Markierungen, die Text visuell hervorheben, ohne das Dokumentenmodell zu verändern -- ein Muster, das dem Ansatz des Prototyps entspricht, jedoch auf dem robusten ProseMirror-Framework aufsetzt.

Das Composable \texttt{useAnalysis} sendet nach einer Debounce-Pause von einer Sekunde den aktuellen Text an das Backend.
Die Antwort wird in Highlighting-Anweisungen übersetzt: Für jedes erkannte Verb wird die Position im Text bestimmt, eine CSS-Klasse zugewiesen (\texttt{verb-good}, \texttt{verb-bad}, \texttt{verb-modal} oder \texttt{verb-neutral}) und ein Tooltip-Text generiert, der die Taxonomiestufe und Quelle der Zuordnung enthält.
Die Highlighting-Extension schreibt diese Anweisungen als \texttt{data-tooltip}-Attribute in die ProseMirror-Decorations.

Bei Textänderungen werden bestehende Decorations automatisch über ProseMirrors Mapping-Mechanismus an die neuen Positionen angepasst, bis die nächste Analyse-Antwort eintrifft.
Dadurch bleiben die Hervorhebungen während des Tippens stabil.

\subsubsection{Tooltip-System}\label{sec:impl-tooltips}

Die Verb-Tooltips ersetzen die nativen Browser-Tooltips des Prototyps durch gestaltete HTML-Elemente.
Ein \texttt{mouseover}-Handler auf dem Editor-Bereich prüft, ob das Ziel-Element ein \texttt{data-tooltip}-Attribut trägt.
Ist dies der Fall, wird ein absolut positioniertes Tooltip-Element oberhalb des Verbs eingeblendet, das die Taxonomiestufe, die Kategorie und -- bei Embedding-basierten Zuordnungen -- die Konfidenz anzeigt.
Die Farbgebung des Tooltips folgt der Verb-Kategorie: Grün für empfohlene, Rot für nicht empfohlene und Blau für Modalverben.

\subsubsection{Score-Bar und Analyse-Panel}\label{sec:impl-scorebar}

Die Score-Bar am unteren Rand des Editors gibt einen kompakten Überblick über die Analyse.
Sie besteht aus drei Sektionen: einem Qualitäts-Score als prozentualer Ringindikator, einer Verb-Statistik mit farbcodierten Zählern (empfohlene, nicht empfohlene und Modalverben) sowie einer Taxonomie-Verteilung, die für jede der sechs Stufen die Anzahl zugeordneter Kompetenzformulierungen anzeigt.

\paragraph{Qualitäts-Score.}

Der Qualitäts-Score aggregiert zwei Dimensionen, die in der Literatur als Qualitätskriterien für Kompetenzformulierungen identifiziert werden.
Die erste Dimension -- \emph{Formulierungsqualität} -- operationalisiert die Forderung der HRK, dass Kompetenzformulierungen \glqq beobachtbares Handeln\grqq\ beschreiben und dazu ein konkretes Handlungsverb verwenden sollen~\cite{hrknexus2015lernergebnisse}.
Sätze mit empfohlenen Verben erhalten den vollen Qualitätswert (1{,}0), Sätze mit mehrdeutigen, aber prinzipiell empfohlenen Verben einen reduzierten Wert (0{,}7), und Sätze mit nicht empfohlenen Verben einen geringen Wert (0{,}2), da die Formulierung zwar als Kompetenz erkannt, aber als zu unspezifisch eingestuft wird.
Nicht als Kompetenzformulierung klassifizierte Sätze tragen nicht zum Score bei.
Die zweite Dimension -- \emph{Taxonomie-Abdeckung} -- folgt der Forderung nach einer breiten Abdeckung kognitiver Anforderungsniveaus~\cite{anderson2001taxonomy}.
Kopf et al.\ empfehlen, dass Modulbeschreibungen verschiedene kognitive Prozessstufen adressieren sollen~\cite{kopf2010kompetenzen}; die Abdeckung misst den Anteil der sechs Stufen, die im Text vertreten sind.
Der Gesamtscore berechnet sich als gewichtete Kombination:
\begin{equation}\label{eq:score}
Q = 0{,}7 \cdot \frac{1}{n} \sum_{i=1}^{n} q(s_i) + 0{,}3 \cdot \frac{|\{\text{belegte Stufen}\}|}{6}
\end{equation}
wobei $q(s_i)$ den Formulierungsqualitätswert des $i$-ten Satzes und $n$ die Gesamtzahl der Sätze bezeichnen.
Die Gewichtung (70/30) ist eine Designentscheidung, die die Formulierungsqualität auf Satzebene gegenüber der Abdeckung priorisiert, da die Korrektur einzelner Formulierungen die primäre Aufgabe des Editors ist.
Die Gewichte wurden nicht empirisch kalibriert; eine aufgabenspezifische Optimierung bleibt zukünftiger Arbeit vorbehalten.

Das seitliche Panel zeigt pro Satz detaillierte Analyseergebnisse: Satztyp (K/I/S), Taxonomiestufe mit Konfidenzwert, Quelle der Zuordnung (Regel oder Embedding) und die fünf ähnlichsten Referenzformulierungen aus der Referenzdatenbank.
Warnungen -- etwa bei Verwendung nicht empfohlener Verben -- werden separat hervorgehoben.

Die ähnlichen Formulierungen werden über alle Sätze dedupliziert, um Redundanz im Panel zu vermeiden, und nach Similarity-Wert sortiert angezeigt.
Zu jeder Empfehlung werden der Similarity-Wert, die Taxonomiestufe und das Herkunftsmodul angegeben, sodass Nutzende die Relevanz der Empfehlung einschätzen können.

\subsubsection{Set-Analyse-Übersicht}\label{sec:impl-set-analyse}

Die Set-Analyse visualisiert die Taxonomie-Verteilung auf Modulebene als Balkendiagramm.
Für jede der sechs Stufen nach Anderson und Krathwohl wird die Anzahl zugeordneter Kompetenzformulierungen dargestellt.
Die Abdeckung -- der Anteil belegter Stufen -- wird numerisch angezeigt.
Fehlende Stufen werden farblich hervorgehoben, und die vom Backend generierte Empfehlung wird als Hinweistext eingeblendet.

\subsubsection{Persistenz und Speicherverwaltung}\label{sec:impl-persistenz}

Die Persistenz nutzt Dexie~4 als IndexedDB-Wrapper.
Das Schema definiert Tabellen für Verblisten (empfohlene, nicht empfohlene und Modalverben), Textbausteine und Kompetenzbeschreibungen.
Die Speicherfunktion unterscheidet -- anders als im Prototyp -- zwischen Neuanlage und Aktualisierung: Bereits gespeicherte Kompetenzbeschreibungen werden per Update aktualisiert, anstatt dupliziert.
Beim Speichern einer noch unbenannten Kompetenzbeschreibung wird ein Dialog eingeblendet, der zur Namenseingabe auffordert; bereits benannte Beschreibungen werden direkt gespeichert.
